\section{GPU-реализация методом Red--Black SOR}

Последовательный метод SOR~\eqref{eq:sor_update} не допускает параллелизма: обновление узла зависит от уже обновлённых соседей на текущей итерации. Архитектура GPU, предоставляющая тысячи параллельных потоков, требует независимости обновлений внутри одного шага. Шахматная (Red-Black) раскраска сетки решает эту проблему: на каждом полуходе обновляемые узлы не зависят друг от друга.

\textit{Шахматная раскраска сетки.}
Каждому узлу~$(i,\,j)$ присваивается цвет в зависимости от чётности суммы индексов:
\begin{equation}\label{eq:rb_coloring}
  c(i,j) = (i + j) \bmod 2.
\end{equation}
Узлы с~$c = 0$ (условно «красные») и~$c = 1$ («чёрные») образуют два непересекающихся подмножества. В пятиточечном шаблоне~\eqref{eq:5point_stencil} каждый красный узел имеет только чёрных соседей и наоборот. Поэтому при одновременном обновлении всех узлов одного цвета каждый из них использует значения соседей противоположного цвета с предыдущего полухода~--- все соседи имеют один и тот же итерационный индекс.

\textit{Обновление давления на полуходе.}
В отличие от классического SOR~\eqref{eq:sor_update}, где часть соседних значений уже обновлена на текущей итерации (индексы~$(m)$ и~$(m+1)$ перемешаны), в Red-Black варианте формула обновления на полуходе цвета~$c$ использует значения всех четырёх соседей с одной и той же итерации:
\begin{equation}\label{eq:rb_sor_step}
\begin{split}
  \tilde{p}_{i,j} &= \frac{1}{a_P}\bigl(
    a_E\,p_{i+1,j}^{(m)} + a_W\,p_{i-1,j}^{(m)}
    + a_N\,p_{i,j+1}^{(m)} + a_S\,p_{i,j-1}^{(m)}
    + b_{i,j}\bigr), \\
  p_{i,j}^{(m+1)} &= p_{i,j}^{(m)}
    + \omega\bigl(\tilde{p}_{i,j} - p_{i,j}^{(m)}\bigr).
\end{split}
\end{equation}
Именно это свойство~--- единый итерационный индекс~$(m)$ у всех соседей~--- обеспечивает независимость обновлений внутри одного цвета и делает метод пригодным для массово-параллельного выполнения на GPU. В ядре обновления дополнительно выполняется проекция~$p \leftarrow \max(p,\, 0)$, соответствующая условию Гюмбеля (подробнее~--- раздел~3.5).

\textit{Граничные условия на GPU.}
После каждой полной итерации (полуход red, затем полуход black) выполняется отдельное ядро, применяющее граничные условия из раздела~3.1: периодичность по окружному направлению~--- связывание ghost-узлов $P[i,\,0] = P[i,\,N_\varphi-2]$, $P[i,\,N_\varphi-1] = P[i,\,1]$; условия Дирихле на торцах~--- $P[0,\,j] = 0$, $P[N_z-1,\,j] = 0$. Граничные условия вынесены в отдельное ядро, чтобы не усложнять логику основного ядра обновления давления.

\textit{Организация вычислений и конфигурация запуска.}
Двумерная сетка потоков GPU формируется из блоков размером~$32 \times 8$: 32~потока по быстрому (окружному) индексу~$j$ соответствуют одной warp-линии, что обеспечивает коалесцированный (последовательный) доступ к глобальной памяти. Каждый поток обрабатывает один внутренний узел; потоки, не удовлетворяющие условию раскраски~$c(i,j) \neq c_{\text{текущий}}$, пропускают обновление.

Коэффициенты дискретизации ($a_E$, $a_W$, $a_N$, $a_S$, $a_P$ и правая часть~$b$) предварительно вычисляются на GPU из поля зазора~$h$ и хранятся в глобальной памяти. В реализации используются обозначения~$A$, $B$, $C$, $D$, $E$, $F$, связанные с обозначениями раздела~3.2 соотношениями $A = a_E$, $B = a_W$, $C = a_N$, $D = a_S$, $E = a_P$, $F = b$.

\textit{Контроль сходимости.}
Сходимость контролируется по критерию относительной поправки~\eqref{eq:convergence_change}. Для снижения накладных расходов проверка выполняется с интервалом в несколько сотен итераций (таблица~\ref{tab:solver_params}): перед блоком итераций сохраняется копия поля давления, после блока вычисляется~$\delta^{(m)}$ средствами массовых операций GPU (поэлементное вычитание, суммирование модулей). При $\delta^{(m)} < \varepsilon$ итерации прекращаются.

\textit{Кэширование GPU-буферов.}
При повторных расчётах на сетке одного размера (например, при параметрическом варьировании эксцентриситета) GPU-буферы под давление и коэффициенты дискретизации переиспользуются. Это устраняет накладные расходы на выделение памяти и повышает производительность серийных расчётов.

\textit{Алгоритм итерации.}
Одна итерация GPU-решателя состоит из трёх этапов: обновление давления в узлах цвета~$c = 0$ (полуход red), обновление в узлах цвета~$c = 1$ (полуход black), применение граничных условий. С заданным интервалом после завершения блока итераций вычисляется относительная поправка~$\delta^{(m)}$ и проверяется критерий остановки. Результирующее поле давления переносится из памяти GPU на CPU.

Описанный GPU-решатель применяется как внутренний решатель линейной системы в алгоритме JFO (раздел~3.5), где дополнительно выполняется обновление активного множества и вычисление степени заполнения~$\vartheta$ в кавитационной зоне.
